\documentclass[12pt]{article}
\usepackage[T1]{fontenc}
\usepackage[utf8]{inputenc}
\usepackage{lmodern}
\usepackage{textcomp}
\usepackage{enumitem}
\usepackage{geometry}
\geometry{margin=1in}
\begin{document}
\begin{center}
Answer Key - Geography - K-12 Level Assessment 14 \
\small (Answers are ordered exactly as in the question paper.)
\end{center}

\section*{Multiple Choice}
\begin{enumerate}[label=\arabic*.]
\item \textbf{Answer:} B
\\ \textit{Options:} A) Abject poverty; B) Better job and higher wages; C) India's caste system; D) Violence in her neighborhood
\\ \textit{Note:} The correct answer is a "pull" factor because better job opportunities and higher wages in the United States attract individuals like Riva to migrate from India in search of improved economic prospects.
\item \textbf{Answer:} C
\\ \textit{Options:} A) Europe; B) East Asia; C) Sub-Saharan Africa; D) South Asia
\\ \textit{Note:} Sub-Saharan Africa is experiencing rapid urbanization due to factors such as population growth, rural-to-urban migration, and economic development, making it the most rapidly urbanizing area in the world.
\item \textbf{Answer:} C
\\ \textit{Options:} A) Ural-Altaic; B) Basque; C) Indo-European; D) Phoenician
\\ \textit{Note:} Most European languages are classified as Indo-European because they share a common ancestral language, which has evolved into the various language families spoken across Europe today.
\item \textbf{Answer:} B
\\ \textit{Options:} A) separatism.; B) gerrymandering.; C) containment.; D) domino theory.
\\ \textit{Note:} Gerrymandering refers to the manipulation of electoral district boundaries to favor a particular political party, thereby undermining fair representation in the voting process.
\item \textbf{Answer:} A
\\ \textit{Options:} A) possibilism.; B) animism.; C) environmental determinism.; D) cultural ecology.
\\ \textit{Note:} Possibilism is the theory in geography that emphasizes human agency and the ability of people to shape their cultural development through choices, rather than being solely determined by environmental factors.
\end{enumerate}

\section*{True / False}
\begin{enumerate}[label=\arabic*.]
\setcounter{enumi}{5}
\item \textbf{Answer:} True
\\ \textit{Options:} A) True; B) False
\\ \textit{Note:} Supranationalism involves countries coming together to form an alliance where they collectively make decisions and policies, thereby willingly sharing aspects of their sovereignty for mutual benefits.
\item \textbf{Answer:} False
\\ \textit{Options:} A) True; B) False
\\ \textit{Note:} The statement is false because Mandarin Chinese has the largest number of native speakers globally, surpassing English in total speaker count.
\item \textbf{Answer:} True
\\ \textit{Options:} A) True; B) False
\\ \textit{Note:} Hasidic Judaism is considered a distinct religious movement within Judaism that emphasizes its own cultural practices and beliefs without significant blending with other religious traditions, thereby exemplifying its unique identity rather than syncretism.
\item \textbf{Answer:} True
\\ \textit{Options:} A) True; B) False
\\ \textit{Note:} Transnational companies (TNCs) operate across multiple countries and are primarily owned and managed by private entities or shareholders rather than being directly controlled by foreign governments.
\item \textbf{Answer:} False
\\ \textit{Options:} A) True; B) False
\\ \textit{Note:} The Third Agricultural Revolution was primarily driven by advancements in technology and agricultural practices rather than the emergence of a rural workforce, which had already been established prior to this period.
\end{enumerate}

\section*{Long Form Response}
\begin{enumerate}[label=\arabic*.]
\setcounter{enumi}{10}
\item \textbf{Answer:} Bulk-gaining industries are those in which the final product weighs more or has a greater volume than the raw materials used in its production. Bottled orange juice is a prime example of this type of industry because the process of juicing oranges and packaging the juice results in a product that is heavier and more valuable than the individual oranges. In production, oranges are harvested, processed, and then transformed into juice, which is bottled for distribution. This process not only adds value but also creates jobs in agriculture, manufacturing, and distribution, significantly impacting the economy. The demand for bottled orange juice also illustrates how bulk-gaining industries can influence local economies by supporting citrus farming and related industries.
\\ \textit{Note:} correct because it accurately describes how bottled orange juice qualifies as a bulk-gaining industry by illustrating the transformation of lighter raw oranges into a heavier, more valuable final product through processing and packaging.
\item \textbf{Answer:} Japan has made significant efforts to preserve the purity of its native language, particularly in the face of increasing English influence. The government has implemented policies to promote the use of Japanese in education and media, while also encouraging the creation of new Japanese terms to replace borrowed English words. These efforts are rooted in a deep cultural pride and a desire to maintain a unique national identity. By protecting their language, the Japanese society fosters a sense of belonging and continuity, which is essential for cultural heritage. However, balancing modernization with linguistic purity can be challenging, as globalization continues to introduce new ideas and vocabulary into daily life.
\\ \textit{Note:} The answer correctly highlights Japan's proactive measures to safeguard its language from English influence through governmental policies and cultural initiatives, reflecting a strong national identity and societal values.
\end{enumerate}

\vfill
\noindent\textit{End of Answer Key}
\end{document}